\begin{table}[t]
\centering
\caption{Exactitud estandarizada (SA, promedio de 50 pliegues) de cuatro modelos en 7 datasets p\'ublicos de esfuerzo, de menos a m\'as proyectos. Uno es sin error y cero es adivinar al azar.}
\label{tab:sa}
\scriptsize
\setlength{\tabcolsep}{3pt}
\resizebox{\ifdim\width>\columnwidth\columnwidth\else\width\fi}{!}{%
\begin{tabular}{lrrrrr}
\toprule
Dataset & $n$ & Mediana & \shortstack{Log-\\lineal} & \shortstack{Lineal\\transf.} & \shortstack{Bosque\\aleatorio} \\
\midrule
Albrecht & 24 & 0.41 & \textbf{0.76} & 0.34 & 0.63 \\
Maxwell & 62 & 0.36 & \textbf{0.57} & 0.51 & 0.56 \\
COCOMO81 & 63 & 0.48 & 0.60 & \textbf{0.79} & 0.64$^{*}$ \\
Desharnais & 81 & 0.34 & \textbf{0.52} & 0.48 & 0.48 \\
NASA93 & 93 & 0.38 & 0.60 & \textbf{0.65} & 0.64$^{*}$ \\
Kitchenham & 145 & 0.45 & 0.59 & \textbf{0.61} & 0.51 \\
China & 499 & 0.37 & \textbf{0.53} & 0.51 & 0.49 \\
\midrule
Mediana de los 7 & & 0.38 & 0.59 & 0.51 & 0.56 \\
\bottomrule
\end{tabular}%
}
\par\smallskip
{\scriptsize\raggedright $^{*}$ El bosque aleatorio supera al log-lineal (MAE por pliegue, prueba de Wilcoxon, $p<0.05$). En negritas, el mejor modelo de cada dataset.\par}
\end{table}
